\subsection{Reinforcement Learning}
\label{subsec:RL}



To improve upon SFT model, we conduct a unified RLVR (Reinforcement Learning with Verifiable Reward) training stage spanning all available environments, targeting terminal usage, office and productivity workflows, software engineering, search, general tool-calling, math, code, STEM, safety, chat, instruction following, long-context QA, inductive and transductive reasoning, structured outputs, and general model usability. For harness-based environments, we construct training data using a diverse collection of harness implementations and interaction formats, improving robustness to variations in execution settings and reducing overfitting to any particular harness design.

We refresh the RL training data using recent collections and perform reward profiling prior to training. For data mixture and curriculum construction, we adopt the Gaussian-based approach introduced in~\cite{nvidia2025nemotron3nanoopen}. Our training procedure largely follows the asynchronous GRPO algorithm with the stability optimizations proposed in~\cite{nvidia2026nemotron3superopen}, while incorporating several improvements to the training infrastructure, detailed in Section~\ref{subsec:infrastructure}. To support training across a large and diverse set of environments, we use a global batch size of 8192, with each sample generating 16 rollouts. Training begins with a maximum generation length of 48K tokens, which is later increased to 64K tokens. % Unlinking Appendix from here until we actually write something there % Please refer to Appendix \ref{appendix:rlvr} for some targeted upgrades to our RLVR environments compared to \supermodel.
% \todo{folks to add new RLVR envs as appendix sections, if they'd like.}